\section{Implementation and Experimental Setup}
\label{sec:setup}

\subsection{Prototype Implementation}
The prototype is implemented by extending the PEPesc proxy program. The proxy entities run at the two sides of the controlled bottleneck link and intercept traffic using transparent proxy rules. The enhanced modules are integrated into the proxy-side control loop. The perception module supports the original rule-based mode and the hybrid GRU mode. The control module supports the original legacy mode and the proposed adaptive mode. This separation allows the original PEPesc behavior and the enhanced behavior to be evaluated under the same topology and traffic conditions.

The GRU model is trained offline from logged proxy observations and then used online for inference. The inference output is not directly applied to packet transmission. Instead, it is first combined with the rule-based baseline and then processed by the decision and execution layers. This deployment choice follows the conservative design principle of learning-assisted systems: the learning component improves perception, while the final control remains bounded and interpretable.

\subsection{Emulation Topology}
The evaluation uses a Mininet-based four-node topology. Node A acts as the sender, node D acts as the receiver, and nodes B and C host the two PEP entities. The B--C link is the controlled high-latency bottleneck link. Unless otherwise specified, the bottleneck link is configured with a one-way propagation delay of 300~ms, resulting in a baseline RTT of about 600~ms. Bandwidth and loss rate are controlled by Linux traffic control to create static, dynamic, and long-duration switching scenarios.

Traffic is generated by iperf. The proxy logs runtime variables including estimated bandwidth, RTT, queue delay, ACK timing, loss information, FEC state, and control decisions. These logs are used both for model training and for post-experiment analysis.

\subsection{Evaluation Scenarios and Metrics}
We evaluate three groups of scenarios. Static scenarios keep bandwidth and loss rate unchanged during each run and are used to verify whether the enhancement preserves steady-state capability. Dynamic scenarios switch bandwidth and loss rate across multiple stages and are used to evaluate recovery after link changes. Long-duration dynamic scenarios repeat a base dynamic pattern multiple times and are used to test whether the system can maintain effective operation over a longer period.

The primary metrics include receiver throughput, recovery time after link switching, throughput deficit area, queue delay, RTT, and effective running time. Receiver throughput reflects useful delivery performance. Recovery time measures how quickly the system returns to a stable throughput level after a link change. Throughput deficit area captures the cumulative performance gap during recovery. The p95 queue delay and p95 RTT measure tail delay behavior. Effective running time measures the duration for which the connection maintains usable operation under long-duration switching.

\subsection{Compared Systems}
The baseline system is the original PEPesc configuration using rule-based perception and legacy control. The enhanced system uses hybrid GRU perception and adaptive control with ACCEL/HOLD/BRAKE decisions, conservative pacing correction, and the FEC reliability floor. Both systems are evaluated under the same topology, traffic generation method, and scenario configurations.
